\section{Висновки}
\label{sec:conclusions}

\begin{enumerate}
\item
Побудовано аналітичну модель ініціювання відгалуження V-подібної
міжфазної тріщини зсуву в її кутовій точці. Сингулярне поле заданої
стаціонарної тріщини використано як зовнішнє поле для малої
розтягувальної зони передруйнування, орієнтованої вздовж бісектриси
одного з матеріалів. Для цієї зони отримано замкнені співвідношення для
її рівноважної довжини $d_i$, розкриття $\delta_i$, роботи сил опору
відриву $W_i$ та зонального енергетичного параметра
$J_i=\dd W_i/\dd d_i$.

\item
Фізична допустимість локальної зони визначається одночасно умовою
стискального контакту берегів V-подібної міжфазної тріщини,
$Cg_2<0$, та умовою розтягувального нормального поля в матеріалі, де
формується зона, $CQ_i>0$. Для досліджених параметрів ці умови одночасно
виконуються лише для одного з матеріалів. Зміна кута зламу межі поділу
та пружного контрасту змінює області допустимості й рівноважні
характеристики зони.

\item
Утворення зони передруйнування послаблює кутову сингулярність, але не
усуває її повністю. Для фізично допустимих інтервалів базової пари
матеріалів отримано $-1/2<\lambda_i<0$, тому залишкове поле є слабшим
за класичну тріщинну сингулярність $r^{-1/2}$, а пов'язаний із ним
локальний конфігураційний потік пружної енергії прямує до нуля при
стягуванні контуру до кутової точки. Отже, окремий скінченний пружний
доданок до локального енергетичного критерію ініціювання для цієї
кутової задачі не потрібний.

\item
За фіксованої геометрії розкриття зони та її зональний енергетичний
параметр однозначно пов'язані через ту саму рівноважну довжину:
$J_i=\chi_i\sigma_i\delta_i$. Тому деформаційний критерій
$\delta_i=\delta_{i\mathrm{c}}$ та енергетичний критерій
$J_i=J_{i\mathrm{c}}$ визначають той самий критичний стан лише за
узгодженого калібрування незалежно заданих матеріальних параметрів,
$J_{i\mathrm{c}}=\chi_i\sigma_i\delta_{i\mathrm{c}}$. Після
незалежного задання одного з цих критичних параметрів модель дає змогу
визначити критичне зовнішнє навантаження ініціювання відгалуження.

\item
Кількісне використання моделі ґрунтується на маломасштабному
припущенні $d_i\ll l$; основною областю застосовності прийнято
$d_i/l<0.1$, тоді як ділянки кривих із $d_i/l$ до приблизно $0.18$
слід трактувати як екстраполяцію локального наближення. Відгалуження з
кутової точки та міжфазне поширення із зовнішніх вершин V-подібної
тріщини є конкурентними локальними механізмами, а не наперед заданими
послідовними стадіями. Локальна структура розв'язку дає можливість
поширити модель на розвинуті V-подібні міжфазні тріщини шляхом заміни
амплітуди зовнішнього поля відповідною амплітудою кутової
сингулярності; до споріднених застосувань належить також задача
розколювання твердого тіла гострим клином.
\end{enumerate}
